\chapter{Registers}
\label{ch:registers}

\section{Register Map}

The registers occupy the first words of the device segment. Every defined register is 16 bits wide and is read and
written a word at a time. The global registers are at offsets \texttt{0x0000}--\texttt{0x000F}; each voice has 16
registers starting at \texttt{0x0010} plus 16 times the voice number.

\begin{table}[H]
	\centering
	\begin{tabular}{lllp{6.2cm}}
		\toprule
		\textbf{Offset}                  & \textbf{Name}              & \textbf{Access} & \textbf{Description}                                             \\
		\midrule
		\texttt{0x0000}                  & \texttt{CONFIG}            & R/W             & Audio enable                                                     \\
		\texttt{0x0001}                  & \texttt{STATUS}            & R               & Version, sequence interrupt, memory overrun                      \\
		\texttt{0x0002}                  & \texttt{STATUSCLR}         & W               & Clears the memory overrun flag                                   \\
		\texttt{0x0003}                  & \texttt{SEQ\_PERIOD}       & R/W             & Output samples per sequencer tick; 0 stops the ticks             \\
		\texttt{0x0004}                  & \texttt{SEQ\_DONE}         & R, W1C          & Sequence completion flags, one per voice                         \\
		\texttt{0x0005}                  & \texttt{SEQ\_INT\_EN}      & R/W             & Sequence completion interrupt enables                            \\
		\texttt{0x0006}                  & \texttt{INSTR\_BASE}       & R/W             & Window address of the instrument table                           \\
		\texttt{0x0007}                  & \texttt{MASTER\_VOL}       & R/W             & Master volume                                                    \\
		\texttt{0x0008}                  & \texttt{FILTER\_CONFIG}    & R/W             & Filter output: low-pass, band-pass or high-pass                  \\
		\texttt{0x0009}                  & \texttt{FILTER\_CUTOFF}    & R/W             & Filter cutoff                                                    \\
		\texttt{0x000A}                  & \texttt{FILTER\_RESONANCE} & R/W             & Filter resonance                                                 \\
		\texttt{0x000B}                  & \texttt{FILTER\_VOL\_PAN}  & R/W             & Filter output volume and pan                                     \\
		\texttt{0x000C}--\texttt{0x000F} & ---                        & Reserved        & Reads zero, writes ignored                                       \\
		\midrule
		\multicolumn{4}{l}{\textit{Voice registers, at the voice base: \texttt{0x0010} + 16 $\times$ voice}}                                               \\
		\midrule
		base + 0                         & \texttt{CONTROL}           & R/W             & Sequencer enable, routing, portamento                            \\
		base + 1                         & \texttt{VOLUME\_PAN}       & R/W             & Volume and pan                                                   \\
		base + 2                         & \texttt{MOD\_ROUTE}        & R/W             & Modulation destination                                           \\
		base + 3                         & \texttt{INSTR}             & R/W             & Instrument selected for the next note                            \\
		base + 4                         & \texttt{FINE}              & R/W             & Fine pitch adjustment                                            \\
		base + 5                         & \texttt{PORTA\_RATE}       & R/W             & Portamento rate                                                  \\
		base + 6                         & \texttt{COMMAND}           & W               & Direct note or stop                                              \\
		base + 7                         & \texttt{SEQ\_PC}           & R/W             & Sequencer program counter; a write restarts the voice's sequence \\
		base + 8--15                     & ---                        & Reserved        & Reads zero, writes ignored                                       \\
		\bottomrule
	\end{tabular}
	\caption{Register Map}
	\label{tab:register-map}
\end{table}

\emph{W1C} means that writing a 1 to a bit clears it and writing a 0 leaves it unchanged. The names are those of this
manual. A voice register for voice $n$ is at \texttt{0x0010} $+ 16n +$ its offset in the voice block.

\section{Global Registers}

\subsection{Configuration Register}
\label{sec:reg-config}

\begin{figure}[H]
	\centering
	\begin{bytefield}[bitwidth=2em]{16}
		\bitheader{0-15} \\
		\bitbox{1}{\footnotesize EN} &
		\bitbox{15}{Reserved}
	\end{bytefield}
	\caption{Configuration Register (\texttt{CONFIG}, Offset \texttt{0x0000})}
	\label{fig:reg-config}
\end{figure}

\begin{table}[H]
	\centering
	\begin{tabular}{llp{8.0cm}}
		\toprule
		\textbf{Bits} & \textbf{Field}     & \textbf{Meaning}                                                                                            \\
		\midrule
		15            & \texttt{AUDIO\_EN} & Audio enable: when 0 both converter channels are forced to digital zero; when 1 the mix is sent. 0 at reset \\
		14--0         & ---                & Reserved: reads zero, writes ignored                                                                        \\
		\bottomrule
	\end{tabular}
	\caption{Configuration Register Fields}
	\label{tab:reg-config-fields}
\end{table}

In Figure~\ref{fig:reg-config}, EN is \texttt{AUDIO\_EN}. \texttt{AUDIO\_EN} gates the output only. While it is 0 the
device runs exactly as it does when it is 1: envelopes advance, sample positions move, modulation and the filter run,
and the sequencers keep time. Clearing and setting it mutes and unmutes without changing any of that.

\subsection{Status and Status Clear Registers}
\label{sec:reg-status}

\begin{figure}[H]
	\centering
	\begin{bytefield}[bitwidth=2em]{16}
		\bitheader{0-15} \\
		\bitbox{10}{Reserved (reads zero)} &
		\bitbox{1}{\footnotesize OVR} &
		\bitbox{1}{\footnotesize IRQ} &
		\bitbox{4}{VERSION}
	\end{bytefield}
	\caption{Status Register (\texttt{STATUS}, Offset \texttt{0x0001})}
	\label{fig:reg-status}
\end{figure}

\begin{table}[H]
	\centering
	\begin{tabular}{llp{8.0cm}}
		\toprule
		\textbf{Bits} & \textbf{Field}        & \textbf{Meaning}                                                                                                                                                                \\
		\midrule
		15--6         & ---                   & Reads zero                                                                                                                                                                      \\
		5             & \texttt{MEM\_OVERRUN} & A memory deadline was missed: a sample read, or the sequencer work of a tick (Section~\ref{sec:deadlines}). Set by the device and held until cleared through \texttt{STATUSCLR} \\
		4             & \texttt{SEQ\_IRQ}     & 1 while \activelow{IRQ3} is asserted: at least one \texttt{SEQ\_DONE} flag is set and enabled in \texttt{SEQ\_INT\_EN}                                                          \\
		3--0          & \texttt{VERSION}      & Device revision, 1                                                                                                                                                              \\
		\bottomrule
	\end{tabular}
	\caption{Status Register Fields}
	\label{tab:reg-status-fields}
\end{table}

In Figure~\ref{fig:reg-status}, OVR is \texttt{MEM\_OVERRUN} and IRQ is \texttt{SEQ\_IRQ}. \texttt{STATUS} is
read-only; a write to its offset is ignored. A write to \texttt{STATUSCLR} (offset \texttt{0x0002}) with bit 5 set
clears \texttt{MEM\_OVERRUN}; the other bits are ignored, and the register reads zero.

\subsection{Sequencer Registers}
\label{sec:reg-sequencer}

\texttt{SEQ\_PERIOD} (offset \texttt{0x0003}) is an unsigned 16-bit count of output samples per sequencer tick.
Zero stops the ticks. A tick of 666 samples, for example, is about 20.84 ms. A write restarts the count towards the
next tick (Chapter~\ref{ch:sequencer}).

\texttt{SEQ\_DONE} (offset \texttt{0x0004}) holds one sticky flag for each voice, in bits 5--0 (bit $n$ belongs to voice $n$): the device sets a voice's
flag when its sequence ends (Section~\ref{sec:seq-end}). Reading the register does not clear a flag. Writing a 1 to a
bit clears that flag; so does a write to the voice's \texttt{SEQ\_PC}. Bits 15--6 read zero.

\texttt{SEQ\_INT\_EN} (offset \texttt{0x0005}) holds one interrupt enable for each voice, in bits 5--0 (bit $n$ for voice $n$). Bits 15--6 read
zero. Disabling an interrupt never clears a flag. \activelow{IRQ3} is asserted whenever any bit of
\texttt{SEQ\_DONE} and the same bit of \texttt{SEQ\_INT\_EN} are both 1 (Section~\ref{sec:interrupts}).

\subsection{Instrument Base Register}

\texttt{INSTR\_BASE} (offset \texttt{0x0006}) holds the window address (\texttt{0x8000}--\texttt{0xFFFF}) of the first word
of the instrument table (Section~\ref{sec:instrument-table}). Instrument $i$ is at \texttt{INSTR\_BASE} $+ 8i$.

\subsection{Master Volume Register}

\texttt{MASTER\_VOL} (offset \texttt{0x0007}) holds an unsigned 8-bit gain in bits 7--0; bits 15--8 read zero. The mixer
multiplies each output channel by the value and divides by 256 (Section~\ref{sec:master-volume}). The reset value is
255.

\subsection{Filter Registers}
\label{sec:reg-filter}

\begin{figure}[H]
	\centering
	\begin{bytefield}[bitwidth=2em]{16}
		\bitheader{0-15} \\
		\bitbox{8}{Volume} &
		\bitbox{8}{Pan}
	\end{bytefield}
	\caption{Filter Volume and Pan Register (\texttt{FILTER\_VOL\_PAN}, Offset \texttt{0x000B})}
	\label{fig:reg-filter-vol-pan}
\end{figure}

\begin{table}[H]
	\centering
	\begin{tabular}{llp{7.5cm}}
		\toprule
		\textbf{Offset} & \textbf{Register}          & \textbf{Meaning}                                                                                           \\
		\midrule
		\texttt{0x0008} & \texttt{FILTER\_CONFIG}    & Bits 1--0 select the output: 00 low-pass, 01 band-pass, 10 high-pass; 11 is reserved. Other bits read zero \\
		\texttt{0x0009} & \texttt{FILTER\_CUTOFF}    & Unsigned 8-bit cutoff coefficient in bits 7--0 (Section~\ref{sec:filter-cutoff})                           \\
		\texttt{0x000A} & \texttt{FILTER\_RESONANCE} & Unsigned 8-bit resonance in bits 7--0 (Section~\ref{sec:filter-resonance})                                 \\
		\texttt{0x000B} & \texttt{FILTER\_VOL\_PAN}  & Filter output volume in bits 15--8, pan in bits 7--0 (Section~\ref{sec:filter-output})                     \\
		\bottomrule
	\end{tabular}
	\caption{Filter Registers}
	\label{tab:reg-filter}
\end{table}

There is no register bit that turns the filter off. A voice reaches the filter or avoids it by its own routing bit
(Section~\ref{sec:reg-control}).

\subsection{Reset State}
\label{sec:reset}

At reset the device is silent: \texttt{AUDIO\_EN} is 0, so the converter receives digital zero until the program sets
it. The sample generator and the memory scheduler run from reset, muted or not.

\begin{table}[H]
	\centering
	\begin{tabular}{ll}
		\toprule
		\textbf{Item}                                                   & \textbf{State at reset} \\
		\midrule
		\texttt{CONFIG}                                                 & 0                       \\
		\texttt{MASTER\_VOL}                                            & 255                     \\
		\texttt{SEQ\_PERIOD}, \texttt{SEQ\_DONE}, \texttt{SEQ\_INT\_EN} & 0                       \\
		\texttt{MEM\_OVERRUN}                                           & 0                       \\
		Each voice's \texttt{CONTROL}                                   & 0                       \\
		Each voice's volume                                             & 0                       \\
		Each voice's pan                                                & 128                     \\
		Envelopes, playback and modulation                              & Inactive                \\
		Each sequencer's \texttt{WAIT} counter and \texttt{END} counter & 0                       \\
		The filter state                                                & Zero                    \\
		Audio memory                                                    & Undefined               \\
		Every other register                                            & Not yet specified       \\
		\bottomrule
	\end{tabular}
	\caption{State at Reset}
	\label{tab:reset}
\end{table}

A program writes every register it relies on before it starts a voice (Section~\ref{sec:startup}).

\section{Voice Registers}
\label{sec:voice-registers}

\subsection{Control Register}
\label{sec:reg-control}

\begin{figure}[H]
	\centering
	\begin{bytefield}[bitwidth=2em]{16}
		\bitheader{0-15} \\
		\bitbox{1}{\footnotesize SEQ} &
		\bitbox{1}{\footnotesize MOD} &
		\bitbox{1}{\footnotesize FLT} &
		\bitbox{1}{\footnotesize POR} &
		\bitbox{12}{Reserved}
	\end{bytefield}
	\caption{Voice Control Register (\texttt{CONTROL}, Voice Offset 0)}
	\label{fig:reg-control}
\end{figure}

\begin{table}[H]
	\centering
	\begin{tabular}{llp{8.0cm}}
		\toprule
		\textbf{Bits} & \textbf{Field}      & \textbf{Meaning}                                                                                                              \\
		\midrule
		15            & \texttt{SEQ\_EN}    & Sequencer enable: the voice's sequencer runs. When 0 the CPU controls the voice directly (Section~\ref{sec:direct-control})   \\
		14            & \texttt{MOD\_MODE}  & Modulation source: the voice makes no sound and steers a parameter (Chapter~\ref{ch:modulation})                              \\
		13            & \texttt{FILTER\_EN} & Filtered: the voice goes to the filter input instead of the dry mix                                                           \\
		12            & \texttt{PORTA\_EN}  & Portamento: a note on the instrument that the voice is already playing glides to the new pitch (Section~\ref{sec:portamento}) \\
		11--0         & ---                 & Reserved                                                                                                                      \\
		\bottomrule
	\end{tabular}
	\caption{Voice Control Register Fields}
	\label{tab:reg-control-fields}
\end{table}

In Figure~\ref{fig:reg-control}, SEQ, MOD, FLT and POR are \texttt{SEQ\_EN}, \texttt{MOD\_MODE}, \texttt{FILTER\_EN}
and \texttt{PORTA\_EN}. \texttt{MOD\_MODE} takes precedence over \texttt{FILTER\_EN}: a voice in modulation mode is in
neither mix. Changing \texttt{SEQ\_EN} does not start, stop or restart anything. A voice that was playing keeps
playing, and a sequence keeps its program counter.

\subsection{Volume and Pan Register}

\begin{figure}[H]
	\centering
	\begin{bytefield}[bitwidth=2em]{16}
		\bitheader{0-15} \\
		\bitbox{8}{Volume} &
		\bitbox{8}{Pan}
	\end{bytefield}
	\caption{Voice Volume and Pan Register (\texttt{VOLUME\_PAN}, Voice Offset 1)}
	\label{fig:reg-volume-pan}
\end{figure}

Bits 15--8 are the voice's volume, unsigned: 0 is silent. Bits 7--0 are its pan: 0 is the left channel only, 128 is
centre and 255 is almost entirely the right channel (Section~\ref{sec:mixing}). The reset volume is 0 and the reset pan
is 128. A volume or pan that a modulation voice is steering changes the value in force, not the value in the register
(Chapter~\ref{ch:modulation}).

\subsection{Modulation Route Register}
\label{sec:reg-mod-route}

\begin{figure}[H]
	\centering
	\begin{bytefield}[bitwidth=2em]{16}
		\bitheader{0-15} \\
		\bitbox{3}{DEST} &
		\bitbox{3}{VOICE} &
		\bitbox{4}{SHIFT} &
		\bitbox{1}{\footnotesize INV} &
		\bitbox{5}{Reserved}
	\end{bytefield}
	\caption{Modulation Route Register (\texttt{MOD\_ROUTE}, Voice Offset 2)}
	\label{fig:reg-mod-route}
\end{figure}

\begin{table}[H]
	\centering
	\begin{tabular}{llp{8cm}}
		\toprule
		\textbf{Bits} & \textbf{Field} & \textbf{Meaning}                                                                                                  \\
		\midrule
		15--13        & \texttt{DEST}  & Destination: 0 off, 1 pitch, 2 volume, 3 pan, 4 filter cutoff, 5 master volume, 6 filter pan; 7 reserved          \\
		12--10        & \texttt{VOICE} & Destination voice, 0--5, for the destinations that belong to a voice (pitch, volume, pan); ignored for the others \\
		9--6          & \texttt{SHIFT} & Signed 4-bit scale exponent, $-8$ to $+7$                                                                         \\
		5             & \texttt{INV}   & Invert the modulation value                                                                                       \\
		4--0          & ---            & Reserved                                                                                                          \\
		\bottomrule
	\end{tabular}
	\caption{Modulation Route Register Fields}
	\label{tab:reg-mod-route-fields}
\end{table}

The register takes effect only while the voice's \texttt{MOD\_MODE} is 1 (Chapter~\ref{ch:modulation}).

\subsection{Instrument and Fine Registers}

\texttt{INSTR} (voice offset 3) holds the number of the instrument the voice's next note uses, in bits 7--0; bits 15--8
read zero. \texttt{FINE} (voice offset 4) holds a signed 13-bit pitch adjustment in bits 12--0. A read of \texttt{FINE}
returns the value sign-extended to 16 bits. Both are set by the sequencer's \texttt{INST} and \texttt{FINE}
instructions, or by the CPU while \texttt{SEQ\_EN} is 0. While \texttt{SEQ\_EN} is 1, a CPU write to either is
discarded.

\subsection{Portamento Rate Register}

\texttt{PORTA\_RATE} (voice offset 5) is an unsigned 16-bit rate: how far the step moves at each envelope update, in raw units of the 16.16 step, where one unit is 1/65,536 of a sample byte per output sample. An envelope update is made every 64 output samples, 2.0 ms (Section~\ref{sec:portamento}). Zero holds the current pitch.

\subsection{Command Register}
\label{sec:reg-command}

\begin{figure}[H]
	\centering
	\begin{bytefield}[bitwidth=2em]{16}
		\bitheader{0-15} \\
		\bitbox{1}{\footnotesize N} &
		\bitbox{8}{Reserved} &
		\bitbox{7}{NOTE}
	\end{bytefield}
	\caption{Voice Command Register (\texttt{COMMAND}, Voice Offset 6)}
	\label{fig:reg-command}
\end{figure}

In Figure~\ref{fig:reg-command}, N is bit 15. \texttt{COMMAND} is write-only and reads zero. With bit 15 set, a write
starts the note in bits 6--0, a number from 0 to 119 (Section~\ref{sec:note-numbers}). With bit 15 clear, a write stops
the voice: its envelope enters release; the sound is not cut (Chapter~\ref{ch:envelope}). The other bits are reserved.
A write is accepted only while \texttt{SEQ\_EN} is 0; with \texttt{SEQ\_EN} set it is ignored.

A note can need the instrument and pitch table to be read from audio memory first, so it can take more than one output
sample to start sounding (Section~\ref{sec:note-start}). The device holds the command until it can apply it. It holds
one command for each voice: a later write replaces a command that has not yet been applied.

\subsection{Sequencer Program Counter}

\texttt{SEQ\_PC} (voice offset 7) holds the window address of the voice's next sequencer instruction, a value in
\texttt{0x8000}--\texttt{0xFFFF}; bit 15 always reads 1. A read returns the current program counter. A write restarts
the voice's sequence (Section~\ref{sec:seq-restart}).
